\documentclass[11pt,letterpaper]{article}
\usepackage[T1]{fontenc}
\usepackage{lmodern,amsmath,amssymb,amsthm,booktabs}
\usepackage[margin=0.94in,top=0.91in,bottom=0.90in,headheight=14pt,headsep=15pt]{geometry}
\usepackage{xcolor,fancyhdr}
\definecolor{ink}{HTML}{173C4A}
\definecolor{accent}{HTML}{146961}
\usepackage[colorlinks=true,linkcolor=accent,urlcolor=accent,citecolor=accent]{hyperref}
\hypersetup{pdftitle={Rigidity and gaps in Euler-factor matching},pdfauthor={Research note prepared for Vladimir Reshetnikov}}
\pagestyle{fancy}\fancyhf{}
\fancyhead[L]{\small\sffamily\color{ink} Euler-factor matching}
\fancyhead[R]{\small\sffamily\color{ink} Dirichlet character packets}
\fancyfoot[C]{\small\sffamily\thepage}
\renewcommand{\headrulewidth}{.25pt}
\setlength{\parindent}{1em}
\setlength{\parskip}{2pt}
\setlength{\emergencystretch}{2em}
\newcommand{\F}{\mathbb F}
\newcommand{\E}{\mathbb E}
\newcommand{\Z}{\mathbb Z}
\newcommand{\C}{\mathbb C}
\newcommand{\ind}{\mathbf 1}
\newcommand{\supp}{\operatorname{supp}}
\newcommand{\ord}{\operatorname{ord}}
\newcommand{\abs}[1]{\lvert#1\rvert}
\newtheorem{theorem}{Theorem}
\newtheorem{lemma}[theorem]{Lemma}
\newtheorem{corollary}[theorem]{Corollary}
\newcommand{\heading}[1]{\section*{\color{ink}#1}}
\newcommand{\subhead}[1]{\par\smallskip\noindent\textbf{#1.}\enspace}
\newcommand{\pin}{c0f2782360ddffd79a20c6d828d582e214a2b8c3}
\newcommand{\gh}[2]{\href{https://github.com/VladimirReshetnikov/ProveIt/blob/c0f2782360ddffd79a20c6d828d582e214a2b8c3/Analysis/GeneralizedRiemannHypothesis/#1}{#2}}

\newcommand{\miniGrid}{\linethickness{0.22pt}\color[gray]{0.79}%
  \multiput(0,0)(1,0){5}{\line(0,1){4}}\multiput(0,0)(0,1){5}{\line(1,0){4}}}
\newcommand{\ptA}[2]{\put(#1,#2){\color{accent}\circle*{.43}}}
\newcommand{\ptB}[2]{\put(#1,#2){\color{white}\circle*{.48}}\put(#1,#2){\color{ink}\linethickness{.9pt}\circle{.48}}}
\newcommand{\miniPlane}[1]{\setlength{\unitlength}{11pt}\begin{picture}(5,5)(-.5,-.5)\miniGrid #1\end{picture}}

\begin{document}
\thispagestyle{plain}
\noindent{\sffamily\small\color{accent} RESEARCH NOTE\hfill 27 NOVEMBER 2026}\par
\vspace{8pt}
\noindent{\color{ink}\LARGE\bfseries Rigidity and gaps in Euler-factor matching}\par
\vspace{3pt}\noindent{\large Dirichlet characters of odd prime order}\par
\vspace{8pt}\noindent\rule{\linewidth}{0.5pt}

\noindent\textbf{Abstract.}
How often must two different products of Dirichlet $L$-functions disagree at a prime?
For characters of exponent $p$, an odd prime, let $m$ be the number of factors on each side after cancellation. The disagreement density is at least $(p-1)/(pm)$. Equality holds exactly for distinct parallel affine subspaces. Outside this family the density is at least $2(1-1/p)$ times the benchmark, sharply for every $p$. Parallelograms and pairs of intersecting lines account for every case at the second level. At packet sizes $k\le p$ we obtain a different, often stronger lower bound and characterize exactly the sizes for which it is sharp. The obstruction comes from an embedded dual $3$-net. Under Dirichlet GRH the sharp coefficient persists in finite prime intervals, with an error costing only a logarithm of $p$. A uniform version retains the factor $4/3$.

\heading{1. From Euler factors to a finite problem}
Fix a modulus $Q$ and an elementary abelian group $X$ of its Dirichlet characters, of rank $r\ge1$ over $\F_p$. Fix a primitive $p$th root $\zeta$ and write $V=X^*$. For a prime $\ell\nmid Q$ there is an evaluation vector $v_\ell\in V$ such that $\chi(\ell)=\zeta^{v_\ell(\chi)}$. A packet is a multiset $A$ of $k$ members of $X$. The reciprocal Euler factor of its product at $\ell$ is
\[
 P_A(\ell,T)=\prod_{\chi\in A}(1-\chi(\ell)T).
\]
For distinct equal-sized packets $A,B$, let $\delta$ denote the natural density within the primes of $P_A(\ell,T)\ne P_B(\ell,T)$. Write $c(x)=\operatorname{mult}_A(x)-\operatorname{mult}_B(x)$ and $m=\sum_x c(x)_+=\tfrac12\sum_x|c(x)|$. Thus $1\le m\le k$. At $v\in V$ the two multisets of exponents differ precisely on the set
\[
 S=\{v:H(v)\ne0\},\qquad H(v)=\sum_{x\in X}c(x)\zeta^{v(x)},\qquad
 \delta=|S|/p^r.                                                    \tag{1}
\]
Indeed, the difference of the projected multiplicities has generating polynomial of degree at most $p-1$, coefficient sum zero, and value $H(v)$ at $\zeta$. The minimal polynomial $1+T+\cdots+T^{p-1}$ implies that a zero value forces every coefficient to vanish. Equality of the traces therefore decides the entire local polynomial. The map $(\Z/Q\Z)^\times\to V$ is onto: the annihilator of its image consists of characters in $X$ trivial on every unit. The prime number theorem in arithmetic progressions now gives (1). The set $S$ is closed under nonzero scalar multiplication and omits the origin.

The finite geometry and density conclusions below are unconditional. The vendored Dirichlet zero-location theorem enters in Section~6, where we pass from density to primes in a specified interval.

\begin{center}
{\small\sffamily\color{ink} Three examples in the character plane $\F_5^2$}\smallskip

\begin{tabular}{ccc}
\miniPlane{\ptA{0}{3}\ptA{1}{3}\ptA{2}{3}\ptA{3}{3}\ptA{4}{3}\ptB{0}{1}\ptB{1}{1}\ptB{2}{1}\ptB{3}{1}\ptB{4}{1}}
&\miniPlane{\ptA{1}{1}\ptA{3}{3}\ptB{1}{3}\ptB{3}{1}}
&\miniPlane{\ptA{0}{2}\ptA{1}{2}\ptA{3}{2}\ptA{4}{2}\ptB{2}{0}\ptB{2}{1}\ptB{2}{3}\ptB{2}{4}}\\[-1pt]
{\small Parallel lines}&{\small Opposite pairs}&{\small Punctured lines}\\
{\small $m=5,\ \delta=4/25$}\quad&\quad{\small $m=2,\ \delta=16/25$}\quad&\quad{\small $m=4,\ \delta=8/25$}
\end{tabular}

{\small Filled and open points give the two packets. Their normalized ratios are $1,\ 8/5,\ 8/5$.}
\end{center}

\newpage
\heading{2. The first two possible disagreement levels}
Put $b_m=(p-1)/(pm)$. An affine flat is a coset of a linear subspace, and distinct cosets of the same subspace are parallel.

\begin{theorem}[Sharp isolation and both equality cases]\label{thm:gap}
For every distinct pair, $\delta\ge b_m$. Equality holds exactly when
\[
             c=\ind_{a+P}-\ind_{b+P},\qquad a-b\notin P              \tag{2}
\]
for a subspace $P\le X$. In particular $m=|P|$. For $r\ge2$ the exact next value of the normalized ratio $\delta/b_m$, over all packet sizes, is
\[
                       T_p=2(1-1/p).                                \tag{3}
\]
Every pair outside (2) has ratio at least $T_p$. Equality holds precisely for the following two configurations after quotient by the period subspace of $c$, up to translation and exchange of the two sides.
\begin{itemize}
\item Opposite vertex pairs of a parallelogram, $\{0,u+w\}$ and $\{u,w\}$ with $u,w$ independent.
\item Two intersecting affine lines, each with their common point deleted.
\end{itemize}
Each point of a quotient configuration represents its full coset with multiplicity one. At $p=3$ the two descriptions coincide and $T_3=4/3$. In rank one the ratio is $m$, and the first two attainable values are $1$ and $2$.
\end{theorem}

In particular the least positive disagreement for $k=p^j$, $0\le j<r$, is $(p-1)/p^{j+1}$. For $k\ge p^{r-1}$ the minimum stays at $(p-1)/p^r$: every nonempty $S$ contains a punctured line, and common factors can pad the hyperplane construction up to $k$.

\subhead{An observable lower bound}
Give $V$ uniform probability. After exchanging signs choose $a$ with $c(a)=\|c\|_\infty$. Define the real-valued function
\[
 g_a(v)=\frac1m\sum_{x\in X}c(x)\ind_{\{v(x-a)=0\}}.
\]
It is the normalized excess of the first packet in the projected fiber containing $a$. It is zero off $S$ and at most one on $S$. For $x\ne a$ the indicator has mean $1/p$, while for $x=a$ its mean is one. Since $\sum c=0$,
\[
        g_a\le\ind_S,\qquad \E g_a=b_m c(a),\qquad
        \delta\ge b_m\|c\|_\infty.                                  \tag{4}
\]
A repeated coefficient after cancellation already gives $\delta\ge 2b_m$.

\subhead{The role of periods}
It remains to treat $c=\ind_A-\ind_B$ with $A,B$ disjoint sets of size $m$. Let $P=\{t:c(x-t)=c(x)\ \forall x\}$. This is a subspace, and each set is a union of $h=m/|P|$ cosets. For $h=1$ we get (2). Its Fourier sum is
\[
 H(v)=|P|\,\ind_{P^\perp}(v)(\zeta^{v(a)}-\zeta^{v(b)}).
\]
The linear form $v\mapsto v(a-b)$ on $P^\perp$ is nonzero. Its kernel has relative size $1/p$, so $\delta=b_m$. We next quantify the cost of a second coset.

\newpage
\heading{3. How a second coset forces a gap}
Fix $a\in A$. On $S$, subtracting the definition of $g_a$ from one and dropping the nonnegative contribution from $B$ gives
\[
 \ind_S-g_a\ \ge\ \frac1m\sum_{y\in A}
          \ind_{S\cap\{v:\ v(y-a)\ne0\}}.                          \tag{5}
\]
Off $S$ both sides vanish. For $y$ outside $a+P$, put $t=y-a$ and $c_t(x)=c(x)-c(x-t)$. This is a nonzero balanced integer function. Its Fourier sum is $H(v)(1-\zeta^{v(t)})$, so its disagreement set is the summand in (5). Write $m_t=\tfrac12\sum_x|c_t(x)|$. The positive list before cancellation is $A$ together with $B+t$, and the negative list is $B$ together with $A+t$. Because $a+P\subset A$ and $y+P\subset A$, the $|P|$ points in $y+P$ cancel between $A$ and $A+t$. Hence
\[
                         m_t\le2m-|P|.
\]
By (4) for $c_t$, the mean of that summand is at least $(1-1/p)/(2m-|P|)$. Exactly $m-|P|$ choices of $y$ are outside $a+P$. Taking means in (5) proves the more precise bound
\[
 \boxed{\displaystyle\quad
  \frac{\delta}{b_m}\ \ge\ 1+\frac{m-|P|}{2m-|P|}
       =\frac{3h-2}{2h-1}.\quad}                                    \tag{6}
\]
For $h\ge2$ the right hand side is at least $4/3$. For $h\ge3$ it is at least $7/5$. Together with (4) this proves the gap and the necessity of (2) for equality at the first level.

\subhead{Differences that attain the first level}
Pass to the quotient by $P$. The signed function is a pullback of its quotient function, and its Fourier sum vanishes off $P^\perp$. Thus $m$ decreases by the factor $|P|$ while $\delta$ increases by that factor. The ratio $\delta/b_m$ is unchanged. Work with period-free sets of size $h$ and $t=y-a\ne0$ as above. If $c_t=\ind_{u+D}-\ind_{v+D}$ with $|D|=N$, then
\[
                            N\le h/(1-1/p).                           \tag{6a}
\]
Indeed, $c_t$ sums to zero along each line parallel to $t$, by telescoping. Consequently $t\notin D$ and the two cosets project bijectively onto the same $N$ points modulo $\langle t\rangle$. It follows that $v+D=u+D+st$ for a single $1\le s\le p-1$. On each of the $N$ lines there is one upward jump and one downward jump, separated by $s$ steps. The values of $c$ are $0,1$ on a line of positive type or $-1,0$ on one of negative type. The former uses $s$ positive points, the latter $p-s$ negative points. If $j$ lines are positive then
\[
      h\ge \max\{js,(N-j)(p-s)\}\ge Ns(p-s)/p\ge N(1-1/p).
\]
This proves (6a). It gives extra information precisely when the general inequality $m_t\le2h-1$ would give too little.

\newpage
\heading{4. An arithmetic condition on packet size}
The preceding theorem is strongest near its first level. For a packet with few factors relative to $p$, matching even one specified constituent gives a different restriction. Define
\[
               D_p(k)=\frac{(p-1)(p+1-k)}{p^2}.
\]
\begin{theorem}[The small-packet bound and its exact sharpness]\label{thm:net}
Let $r\ge2$ and $1\le k\le p$. Every distinct pair of size $k$ has $\delta\ge D_p(k)$. For each such rank, a pair with equality exists if and only if
\[
                    k\mid(p-1)\qquad\text{or}\qquad k\mid(p+1)\qquad\text{or}\qquad k=p.
                                                                    \tag{7}
\]
If $k$ satisfies none of these conditions, then
\[
                   \delta\ \ge\ \frac{p+1-k}{p}+\frac{k-2}{p^r}.     \tag{8}
\]
For $r=1$, every distinct equal-sized pair has $\delta=1-1/p$.
\end{theorem}

\subhead{Counting hyperplanes}
Cancel common factors. Fix $a$ in the first remaining packet. Agreement forces $v(a)=v(b)$ for at least one member $b$ of the other packet. Thus the agreement set lies in the union of at most $m$ hyperplanes $(a-b)^\perp$. The first has $p^{r-1}$ points. Every further one adds at most $p^{r-1}-p^{r-2}$, by intersection with the first. It follows that $\delta\ge D_p(m)\ge D_p(k)$. In rank one every nonzero projection is injective. This is the point-and-direction argument from discrete tomography \cite{AL}. Lev's theorem on perfect directions also gives the planar bound directly \cite{Lev}.

\subhead{The planar equality problem}
In rank two, write $d$ for the number of agreeing projective directions in $V$. Every line contributes its $p-1$ nonzero vectors, so
\[
                  \delta=\frac{(p-1)(p+1-d)}{p^2},\qquad d\le m\le k. \tag{9}
\]
Represent an evaluation direction by its one-dimensional kernel in $X$. If $d=k$, for every $a$ the $k$ good directions must be the $k$ different directions of the segments $ab$ as $b$ ranges over the other packet. Interchanging the sides proves both packets are disjoint sets and every point has exactly one opposite point in every good direction. Every line joining opposite points contains just one point of either set. Appending the $k$ good directions as points at infinity gives three disjoint sets of size $k$: any line meeting two meets all three once. This is an embedded dual $3$-net.

\subhead{Divisibility from the conic theorem}
For $k\ge3$, Theorem~5.1 of Blokhuis--Korchm\'aros--Mazzocca applies, since $k\le p$ and the third component is collinear. The other two components lie on a possibly reducible conic \cite{BKM}. Their Section~5 describes the resulting chord groups. We explain the divisibility argument and construct every attainable case on the next page. For $k=1,2$, divisibility in (7) is automatic since $p$ is odd.

\newpage
\heading{Conics, constructions and passage to higher rank}
\subhead{Why a conic appears}
Here is the moment argument behind the conic theorem in this setting \cite{BKM}. For each $j<k$, the difference of the two sums of $v(x)^j$ is a homogeneous binary form in $v$ of degree $j$ that vanishes at $k$ projective directions. It is identically zero. Since $j<p$, all the binomial coefficients in its expansion are invertible. The two packets therefore give the same sum for every polynomial on the plane of degree at most $k-1$.

Choose a nonzero quadratic $q$ through three points of $A$ and two of $B$: these are five homogeneous linear constraints on six coefficients. If some $b_0\in B$ lies outside $q=0$, pair the remaining $k-3$ points of $A$ with the remaining $k-3$ points of $B$ other than $b_0$. Multiply $q$ by the equations of their joining lines. Each mixed line contains exactly one point of $B$, so the resulting polynomial vanishes at all packet points except $b_0$. Its degree is at most $k-1$, a contradiction. Thus $q$ contains all of $B$. Use three of those points and two of the original points of $A$ to reverse the argument with the same $q$. This puts both packets on the conic.

\subhead{Which group supplies the directions}
If the conic is a pair of lines, a line containing points of both colors contains just two packet points. For $k\ge3$ one component line must carry at least three points of one color, forcing all points of the other color onto the other line. Each line is therefore monochromatic. A repeated line is impossible, and conjugate nonrational lines contain at most one rational point between them. If the lines intersect in the affine plane, choose them as the axes. Opposite points $(a,0),(0,b)$ determine slopes $-b/a$. For parallel lines choose $(0,a),(1,b)$, giving slopes $b-a$.

For a nonsingular conic, the line at infinity meets it in two rational points, one double rational point, or a conjugate pair. The affine models and their chord parameters are $xy=1$ with parameter $-1/(ab)$, $y=x^2$ with parameter $a+b$, and the norm-one group of $\F_{p^2}^{\times}$ with parameter $-1/(ab)$, respectively \cite[Section~5]{BKM}. In the last model a nonzero direction $z$ is represented by $z^{p-1}$ and $(a-b)^{p-1}=(a^{-1}-b^{-1})/(a-b)$. Each case produces two $k$-subsets $U,W$ of a group of order $p-1$, $p$, or $p+1$, with product set (or sum or difference set) of size $k$. Every $uW$ then equals $UW$. Hence $u_0^{-1}U$ lies in the stabilizer of $W$, which has size at most $|W|=k$, since it acts freely. The stabilizer has size $k$. This proves (7) is necessary in rank two.

\subhead{Three constructions}
If $k\mid p-1$, take $H\le\F_p^\times$ of order $k$ and the packets $H\times\{0\}$ and $\{0\}\times H$. For every slope $-s$, $s\in H$, the pairing $a\mapsto sa$ matches points along that direction. If $k\mid p+1$ and $k\le p$, take a subgroup $H$ of size $k$ in the norm-one group, and a distinct coset $tH$. For each $s\in tH$, the pairing $a\mapsto s/a$ yields the common chord parameter $-1/s$. Finally, if $k=p$, take the two parallel lines $y=1$ and $y=2$. Each direction transverse to the lines produces a bijection between them. In each construction there are $k$ good directions, so (9) is exact. Embedding the plane in $X$ preserves $\delta$ by uniform restriction of $v$ to the plane. If a packet model includes a trivial character and every factor is required to have exact order $p$, translate both packets to avoid zero. There are at most $2k<p^2$ forbidden translations in the plane, and a common translation preserves agreement.

\subhead{A stronger bound when equality is impossible}
For $k$ inadmissible, choose a bad projective point $z$ in $\mathbb P(V)$, which exists by (1) and Fourier inversion. Each projective line through $z$ is a two-dimensional space of evaluations, dual to a quotient of $X$. The projected packets there are distinct since they disagree at $z$, and by the planar result they have at most $k-1$ good projective points. Besides $z$, the line has at least $p+1-k$ bad points. The $(p^{r-1}-1)/(p-1)$ lines through $z$ partition the other points. Multiplication of the resulting count by $(p-1)/p^r$ proves (8), strictly exceeding $D_p(k)$. This completes Theorem~\ref{thm:net} in every rank.

\newpage
\heading{Proof of the sharp second level}
Let $\alpha$ be the infimum of $\delta/b_m$ for all signed packet differences outside (2) at the fixed prime $p$, allowing every rank. A parallelogram has Fourier sum $(1-\zeta^{v(u)})(1-\zeta^{v(w)})$ up to a phase. Its density is $(1-1/p)^2$ and $m=2$, giving $\alpha\le T_p$. Repeated coefficients have ratio at least two. For a simple pair quotient by its periods as in Section~3. The resulting mass $h$ is at least two. In rank one the ratio equals $h$. In higher rank, for $2\le h\le p-1$, the hyperplane count gives
\[
        \delta/b_h\ge h(p+1-h)/p\ge T_p ,
\]
with the quadratic taking its minimum precisely at the endpoints.

Suppose $h\ge p$. In (5), for each $y\ne a$ put $t=y-a$ and let $N$ denote the positive mass of $c_t$. If $c_t$ is outside (2), its density is at least $\alpha b_N$ by definition, and $N\le2h-1$. If it is of form (2), its density is $b_N$ and (6a) gives $h/N\ge1-1/p$. Averaging (5) therefore yields
\[
 \frac{\delta}{b_h}\ge
 \min\!\left\{1+(1-1/h)(1-1/p),\quad 1+\alpha\frac{h-1}{2h-1}\right\}. \tag{9a}
\]
The first entry exceeds $T_p$ by at least $1/p^2$. The second exceeds $\alpha$ by at least $1-T_p p/(2p-1)=1/(2p-1)$. If $\alpha<T_p$, these estimates and the small-$h$ bound put every ratio a fixed positive distance above its infimum. Hence $\alpha=T_p$ and equality is impossible at $h\ge p$.

\subhead{Identifying the two configurations}
Only $h=2$ or $h=p-1$ can attain $T_p$. At $h=2$ the two vector sums must agree, since a nonzero difference of sums forces disagreement of density at least $1-1/p>(1-1/p)^2$. After translation the pairs are $\{0,u+w\},\{u,w\}$. Proportional $u,w$ give density $1-1/p$, whereas independent ones give the required $(1-1/p)^2$.

For $h=p-1\ge4$, equality in the hyperplane count forces the support into an affine plane. To see this, for fixed $a$ the $h$ distinct hyperplanes $(a-b)^\perp$ have union equal to the whole agreement set. Equality in the successive union bounds makes their intersections with each preceding hyperplane equal inside each new hyperplane. In particular all contain the common codimension-two intersection $K$ of the first two. Every hyperplane contained in the union is one of its members: at most $p$ proper linear subspaces cannot cover a vector space, as counting their common origin once shows. Repeating the count at every $a$ now places every difference $a-b$ in $K^\perp$. Translating by a fixed point places both packets in that two-dimensional space.

There are $2p-2>p+1$ points of the two colors there, too many for a nonsingular conic. The classification in Section~4 forces two intersecting lines, since parallel lines would require an additive subgroup of order $p-1$. The intersection is in neither packet, and all other points of each line are present. Conversely the signed difference equals the difference of the \emph{full} line indicators. Its Fourier support in the plane is the union of their two punctured annihilator lines, of density $2(p-1)/p^2$. With $m=p-1$ this gives $T_p$. At $p=3$, two distinct pairs with the same midpoint give both descriptions. This proves Theorem~\ref{thm:gap}.

\newpage
\heading{5. An exact spectrum at exponent seven}
Take $p=7$ and $r=2$. Formula (9) makes the maximum number of good projective directions the only unknown needed for the minimum density. All sizes through seven except $k=5$ meet (7). At five, (8) gives $24/49$ and padding a four-factor extremizer attains it. Above seven, a pair of parallel lines already attains the least possible one-line disagreement. Consequently the following table is exact for \emph{every} packet size.
\begin{center}
\begin{tabular}{lrrrrrrr}
\toprule
Packet size $k$ & 1 & 2 & 3 & 4 & 5 & 6 & $\ge7$\\
\midrule
Least prime disagreement $\times 49$ &42&36&30&24&24&12&6\\
Greatest prime agreement $\times 49$ &7&13&19&25&25&37&43\\
\bottomrule
\end{tabular}
\end{center}
In fact the \emph{entire} set of positive disagreement densities for five factors in rank two is
\[
                \left\{\frac{24}{49},\frac{30}{49},\frac{36}{49},
                               \frac{42}{49},\frac{48}{49}\right\}.   \tag{10}
\]
To see this, each $d=1,2,3,4$ is realized at size $d$ and preserved by padding. To get $d=0$, take two copies of $0$ against independent points $e_1,e_2$. Equality forces $v=0$. Pad to five in the same way. Equation (9) and the impossibility at $d=5$ exhaust the possibilities. A common translation gives exact-order-seven factors if required.

\subhead{An example over the rational primes}
Let $Q=29\cdot43=1247$ and take logarithms to base $2$ modulo $29$ and base $3$ modulo $43$, each reduced modulo seven. They define $v_\ell=(s,t)\in\F_7^2$ on the units. Use the four characters indexed by each of the sets
\[
   A=\{(1,0),(-1,0),(0,1),(0,-1)\},\qquad
   B=\{(2,2),(2,-2),(-2,2),(-2,-2)\}.
\]
A character $(a,b)$ has value $\zeta^{as+bt}$ at such a prime. In $\F_{49}=\F_7[i]$, $i^2=-1$, these sets are the group $\{1,-1,i,-i\}$ and its coset $(2+2i)A$ on the norm-one conic. Precisely the directions $(1,2),(1,3),(1,4),(1,5)$ agree. Including the zero evaluation gives $25$ agreements and $24$ disagreements among $49$ evaluations. Each evaluation occurs in $24$ reduced classes modulo $1247$, so the exact density follows from $576$ distinguishing classes among $1176$. The enumerated primes through $10^6$, omitting $29$ and $43$, gave $38{,}440$ disagreements and $40{,}056$ agreements. The density proof uses the residue classes.

\subhead{An infinite family of packet sizes with the same minimum}
Let $d=p^j$ with $j<r$. For every integer $d\le k\le T_p d$ the minimum at size $k$ is $(1-1/p)/d$. Indeed, (2) gives this value by padding. Every other flat pair has reduced mass a power of $p$ at most $d$, since $T_pd<pd$. Every nonflat pair has $\delta\ge T_p(1-1/p)/k\ge(1-1/p)/d$. Thus the minimum can remain fixed through a whole interval of sizes even as common factors become optional.

\newpage
\heading{6. Distinguishing primes in finite intervals}
Fix a smooth nonnegative $w$ compactly supported in $(1/2,1)$ with integral one. Put
\[
   M_Y(f)=\sum_{\ell\nmid Q\ {\rm prime}} (\log\ell) w(\ell/Y)f(v_\ell).
\]
For $f(v)=\sum_{x\in X}\widehat f(x)\zeta^{v(x)}$ write $\|f\|_{\mathcal A}=\sum_x|\widehat f(x)|$. The article's prime Gram estimate gives, under Dirichlet GRH,
\[
   \left|M_Y(f)-Y\E f\right|\le C_w\|f\|_{\mathcal A}\sqrt{Y}\log(2Q)
                \qquad (Y\ge Y_w).                                  \tag{11}
\]
For one primitive character the smoothed explicit formula weights each nontrivial zero $\rho$ by $Y^\rho\widetilde w(\rho)$, where $\widetilde w$ is the Mellin transform. On the critical line these terms total $O_w(\sqrt{Y}\log(2Q))$ by rapid decay and the local zero count. Removing higher prime powers costs $O_w(\sqrt{Y})$ by Chebyshev's bound, and omitting primes dividing $Q$ costs $O_w(\log(2Q))$. Summing the single-character estimate proves (11). See \gh{Article/sections/rec.tex}{the article's prime Gram lemma} for the contour argument.

\begin{theorem}[A finite form of the sharp gap]\label{thm:finite}
Let $c$ be the reduced difference and let $\Lambda=1$ for the two-coset case (2), $\Lambda=4/3$ otherwise. Under Dirichlet GRH,
\[
    M_Y(\ind_S)\ \ge\ \Lambda b_m Y-O_w(\sqrt{Y}\log(2Q)).             \tag{12}
\]
For (2) the estimate is two-sided. The error constant is independent of $p,r,k,m,Q$. Each interval $(Y/2,Y)$ contains a distinguishing prime once $Y\ge C'_w m^2\log^2(2Q)$. For simple nonflat packets, $(3h-2)/(2h-1)$ can replace $4/3$. For repeated coefficients $\|c\|_\infty$ can replace it.
\end{theorem}

\begin{proof}
For $x\in X$, $\ind_{v(x-a)=0}=p^{-1}\sum_{t\in\F_p}\zeta^{v(t(x-a))}$. In the formula for $g_a$ the constant $t=0$ terms cancel. The triangle inequality gives $\|g_a\|_{\mathcal A}\le2(1-1/p)$. It is a pointwise minorant of $\ind_S$ with the mean in (4). In the two-coset case, it equals $\ind_S$ because their means agree.

For simple nonflat packets, strengthen the minorant using (5). For each of the $m-|P|$ terms indexed by $y\notin a+P$ choose an anchor and overall sign for $c_{y-a}$ that makes its maximal absolute coefficient positive. Its own $g$ is a minorant of its active-set indicator, has norm at most $2(1-1/p)$ and mean at least $(1-1/p)/(2m-|P|)$. Add $1/m$ times the sum of these minorants to $g_a$. The resulting function lies below $\ind_S$, has norm at most $4(1-1/p)$, and its mean is at least the right hand side of (6) times $b_m$. Positivity of $w$ and (11) give (12). For a repeated coefficient use (4) directly. Finally $b_m\ge2/(3m)$, so a sufficiently large $C'_w$ ensures positive weight in the interval.
\end{proof}

For an exact second-level pair there is also a two-sided estimate $M_Y(\ind_S)=T_p b_mY+O_w(\sqrt Y\log(2Q))$ uniformly in $p$. For a parallelogram on the period quotient the distinguishing indicator is the product of the indicators of two nonzero linear evaluations. Each factor has Fourier norm at most two. For punctured lines it is $\ind_H+\ind_{H'}-2\ind_{H\cap H'}$ on the dual of the quotient, where $H,H'$ annihilate the two line directions. Subspace indicators have norm one, so again the norm is at most four. Extension by zero off $P^\perp$ preserves this norm.

\newpage
\heading{A finite interval with the sharp prime-dependent coefficient}
For nonflat reduced packets there is also a minorant of mean at least $T b_m$, where $s=1-1/p$ and $T=2s$, with Fourier norm $O(\log(2p))$. Here $b_m=s/m$. Hence
\[
 M_Y(\ind_S)\ge T b_mY-O_w\!\left(\log(2p)\sqrt Y\log(2Q)\right).
\]
We give the construction because a density estimate alone need not have a bounded-norm witness. If a coefficient has absolute value at least two, (4) suffices. Otherwise pass to the quotient by the period group, with $n$ points of each sign. Extending any quotient witness by zero outside $P^\perp$ divides its mean by $|P|$ and preserves its Fourier norm: each coefficient is divided by $|P|$ and repeated $|P|$ times.

For $2\le n\le(p-1)/2$ enumerate the $n$ anchor hyperplanes $H_j=(a-b_j)^\perp$. The real minorant
\[
 F=(1-\ind_{H_1})\left(1-\sum_{j=2}^n\ind_{H_j}\right)
\]
has mean at least $D_p(n)$ by the hyperplane-intersection count, and norm at most $2n$. Scale it by $T b_n/D_p(n)\le1$, using $n(p+1-n)/p\ge T$. The norm is now at most $2Tp/(p+1-n)\le8$.

For larger $n$ use successive minorants in (5), starting at the anchor guarantee $L_0 b_n$ with $L_0=1$. If $c_t$ is flat, the cycle bound gives $m_t\le n/s$. For $n\ge p-1$ this implies $1+(n-1)/m_t\ge T$; for $(p+1)/2\le n\le p-2$, integrality gives $m_t\le n$ and again $1+(n-1)/m_t\ge T$. Its exact flat indicator has norm at most two, as follows from the two-coset case of (4). For nonflat $c_t$ the preceding iteration gives mean at least $L_i s/m_t$ and $m_t\le2n-1$. Combining the $n-1$ terms of (5) with the anchor, all nonflat pairs admit a new minorant of mean at least $L_{i+1}b_n$, where
\[
 L_{i+1}=\min\{T,1+\alpha L_i\},\qquad
 \alpha=\frac{p-1}{2p}\le\frac{n-1}{2n-1}.
\]
Indeed a mixture of the flat and nonflat choices is bounded below by the smaller of the two corresponding pure guarantees. At each iteration the norm bound increases by at most two, with the small-$n$ bound of eight available at every stage. The affine recursion has fixed point $2p/(p+1)=T+2/[p(p+1)]$ and contracts by less than $1/2$. Thus $J=\lceil\log_2(p(p+1))\rceil$ steps suffice, with norm at most $8+2J$. Applying (11) yields the claim.

\subhead{What the two finite bounds measure}
Equation (12) gives a distinguishing prime at a scale uniform in $p$. The estimate on this page recovers the full coefficient $2(1-1/p)$ at each fixed exponent. For any $0<\varepsilon<1$, the smoothed distinguishing-prime mass is at least $(T-\varepsilon)b_mY$ once $Y\ge C_w\varepsilon^{-2}m^2\log^2(2p)\log^2(2Q)$. The bound counts the primes themselves with nonnegative weights and applies in every such dyadic interval.

\newpage
\heading{7. Sources, verification and remaining questions}
\subhead{The two repository inputs}
We used ProveIt commit \texttt{c0f2782360dd}. Its \gh{vendor/grh-lean/README.md}{vendor manifest} identifies upstream commit \texttt{72fead102f66} in \href{https://github.com/lsernau/grh-lean/tree/72fead102f663d5c7a407964a13dc0cf9d977207}{\texttt{lsernau/grh-lean}}. The theorem \texttt{LFunction\_eq\_zero\_re} in \gh{vendor/grh-lean/grh/src/Main.lean}{\texttt{Main.lean}} has the \gh{vendor/grh-lean/grh/src/Statements.lean}{interface} that every zero in $0<\Re s<1$ has $\Re s=1/2$, for all Dirichlet characters of every positive modulus. This is the analytic input in (11). The upstream \gh{vendor/grh-lean/grh/README.md}{verification notes} record the pinned Lean and Mathlib versions and independent export checks. The finite-interval argument uses this declaration as an input. The reported upstream checks were not rerun here. The present extension has a mathematical proof and finite executable checks.

The \gh{Article/README.md}{merged article} treats packets of quadratic characters in \gh{Article/sections/rec.tex}{the section on recovering zeros and characters}. Its theorem labelled \texttt{rec:thm:qmatch} determines a high-agreement range for exponent two. The end of that section explicitly asks for an odd-prime analogue. Theorems~\ref{thm:gap} and~\ref{thm:net} give the first isolated levels and sharpness classification above. Formula (10) supplies a complete spectrum in one instance. The finite estimate uses the article's analytic method and a bounded-norm minorant tailored to the new extremizers.

\subhead{Relationship to published work}
Meshulam's uncertainty inequality already gives the baseline numerical bound for a simple signed support of size $2p^j$ \cite{Mesh}. Finite Radon and Fourier methods for hyperplane sums were developed in \cite{ILMP}. Bir\'o and Lev obtained structural uncertainty bounds in rank two \cite{BL}. In particular, Lev bounds the number of perfect directions by half the support size, apart from a constant function on one line \cite{Lev}. For a balanced nonzero packet difference this gives the planar $d\le m$. Examples with two lines and with parallel cosets already appear in \cite{BL}. The conic and subgroup classification used in Section~4 is due to \cite{BKM}. General Frobenius trace-matching problems for irreducible representations are studied by Walji \cite{Walji}. The contribution proposed here is the all-rank, integer-mass gap with both equality classifications, the sharpness and nonattainment consequences for Euler factors, and their finite-prime form. The literature comparison is targeted and does not establish priority.

\subhead{Finite checks}
The included standard-library Python scripts use integer residues and exact multiplicities. The first enumerates every unordered distinct pair of multisets in $\F_3^2$ of sizes $1,2,3,4$, totalling $136{,}821$ pairs. It also checks all $52{,}650$ pairs of size two in $\F_5^2$ and tests $103$ constructions for primes through $31$. An independent normalized enumeration checks the second equality cases at $p=3,5,7,11$. A further check reconstructs the differences in (6a) for $24$ rank-three examples, including sharp ones. The arithmetic script checks every unit class and the prime sample of Section~5. Their text outputs are included for reproducibility. The arguments in Sections~1--6 cover arbitrary rank and size independently of these computations.

\subhead{Further research questions}
What is the next attainable normalized level above $2(1-1/p)$? Which inadmissible pairs $(p,k)$ have an exact minimum strictly above the planar value in (8) at rank two? How does the actual minimum change beyond rank two at these sizes? Determining the entire odd-prime spectrum for arbitrary $p,r,k$ remains open here.

\newpage
\begin{thebibliography}{9}
\bibitem{AL} A. Alpers and D. G. Larman, \emph{The smallest sets of points not determined by their X-rays}, Bull. Lond. Math. Soc. 47 (2015), 171--176. \href{https://doi.org/10.1112/blms/bdu111}{doi:10.1112/blms/bdu111}.
\bibitem{BL} A. Bir\'o and V. F. Lev, \emph{Uncertainty in finite planes}, J. Funct. Anal. 281 (2021), 109026. \href{https://arxiv.org/abs/1808.07424}{arXiv:1808.07424}.
\bibitem{BKM} A. Blokhuis, G. Korchm\'aros and F. Mazzocca, \emph{On the structure of $3$-nets embedded in a projective plane}, J. Combin. Theory Ser. A 118 (2011), 1228--1238. Theorem 5.1 and Section 5 of the \href{https://arxiv.org/abs/0911.4100}{preprint}. \href{https://doi.org/10.1016/j.jcta.2010.11.007}{doi:10.1016/j.jcta.2010.11.007}.
\bibitem{ILMP} A. Iosevich, A. Liu, A. Mayeli and J. Pakianathan, \emph{Wavelet decomposition and bandwidth of functions defined on vector spaces over finite fields}. \href{https://arxiv.org/abs/1601.03473}{arXiv:1601.03473}.
\bibitem{Lev} V. F. Lev, \emph{Point distribution and perfect directions in $\F_p^2$}, Unif. Distrib. Theory (2020). \href{https://doi.org/10.2478/udt-2020-0012}{doi:10.2478/udt-2020-0012}.
\bibitem{Mesh} R. Meshulam, \emph{An uncertainty inequality for finite abelian groups}, European J. Combin. 27 (2006), 63--67. \href{https://arxiv.org/abs/math/0312407}{arXiv:math/0312407}.
\bibitem{Walji} N. Walji, \emph{Matching densities for Galois representations}, Proc. Amer. Math. Soc. \href{https://doi.org/10.1090/proc/12996}{doi:10.1090/proc/12996}. \href{https://arxiv.org/abs/1501.07343}{arXiv:1501.07343}.
\end{thebibliography}
\end{document}
